% Main-text predictions table for Section 6.4; updating-process rows follow Figure 4 panels A--D.
\begin{table}[tbp]
\centering
\footnotesize
\singlespacing
\setlength{\tabcolsep}{5pt}
\renewcommand{\arraystretch}{1.15}
\setlength{\belowcaptionskip}{6pt}
\caption{Discriminating predictions of the receiver-centered framework. Rows group the three receiver demands (\secref{sec:account-demands}) and the four updating processes of \figref{fig:temporal}A--D. $p_t$ is the receiver's current prediction that danger will follow a sound; $C$ is the cost of an action given current conditions and recent history. No single prediction identifies a mechanism; convergent tests in the same animals are needed. The evidence column lists motivating studies and whether the prediction itself has been tested.}
\label{tab:predictions}
\begin{tabular}{@{}>{\raggedright\arraybackslash}p{1.05in}>{\raggedright\arraybackslash}p{0.9in}>{\raggedright\arraybackslash}p{2.05in}>{\raggedright\arraybackslash}p{2.05in}@{}}
\toprule
\textbf{Process} & \textbf{Where it acts} & \textbf{Discriminating prediction} & \textbf{Current evidence} \\
\midrule
\multicolumn{4}{@{}l}{\textit{Receiver demands}} \\
Rapid acoustic defense
  & Outside the expression (fast route)
  & Present in naïve animals and scales with level, onset, or looming. Hypothesis: transfers across callers when acoustic features are matched.
  & Synthetic sounds in laboratory rodents \citep{DavisEtAl1982,XiongEtAl2015,LiEtAl2021}; untested with natural alarms or acoustically controlled callers. \\
Learned threat evaluation
  & $p_t$
  & Depends on aversive or predator history; specific to call type or caller.
  & Experience-gated freezing to natural calls in rats; caller discounting in juvenile ground squirrels \citep{KimEtAl2010,HareAtkins2001}. \\
Flexible action selection
  & $C$ and the choice among actions
  & The action changes with refuge, food, or state while the sound, the receiver's history, and independent signs of recognition stay the same.
  & Context alters actions in mice and chipmunks \citep{MongeauEtAl2003,BaackSwitzer2000}; untested with an independent recognition measure. \\
\midrule
\multicolumn{4}{@{}l}{\textit{Updating processes}} \\
Habituation
  & History-sensitive responses, including fast ones
  & Partial recovery with rest alone; specific to the repeated stimulus, with transfer to other exemplars or callers in proportion to acoustic similarity.
  & Criteria defined \citep{RankinEtAl2009}; caller-specific decline and renewal in sciurids \citep{Hare1998,BlumsteinDaniel2004} fit habituation but do not identify the mechanism. \\
Reliability learning
  & $p_t$, by source
  & Decline specific to one source that reverses when its calls are again paired with danger, with acoustics and schedule held constant. The direction of change may differ between species.
  & Manipulated reliability changes responding in opposite directions in two sciurids \citep{HareAtkins2001,BlumsteinEtAl2004}; reversal untested. \\
Temporal expectation
  & Expected timing, beyond a single $p_t$
  & Slower decline with irregular than with regular intervals when mean interstimulus interval, stimulus properties, and number of exposures are matched.
  & Artificial sounds only \citep{Davis1970,HerryEtAl2007}; untested with natural calls. \\
Motivational and contextual relevance
  & $C$
  & Response magnitude or choice of action changes when the costs of responding change, with the sound held constant; recovery is one possible outcome.
  & Related observations in laboratory rodents \citep{MongeauEtAl2003,BerzEtAl2022}; untested with costs manipulated and recognition measured. \\
\bottomrule
\end{tabular}
\end{table}
